\documentclass[11pt,a4paper]{article}
\usepackage[margin=2.4cm]{geometry}
\usepackage{amsmath,amssymb,booktabs,graphicx,xcolor,microtype}
\usepackage[colorlinks=true,linkcolor=black,urlcolor=blue!60!black,citecolor=black]{hyperref}
\usepackage{caption}
\captionsetup{font=small,labelfont=bf}
\newcommand{\ok}{\textcolor{green!45!black}{\textbf{+}}}
\setlength{\parskip}{2pt}

\title{\vspace{-1.4cm}\Large Autonomous discovery of symmetry-protected topological phases:\\
benchmark results}
\author{\normalsize Yuxuan Zhang}
\date{\normalsize 23 August 2026}

\begin{document}
\maketitle
\vspace{-0.7cm}

\begin{abstract}\noindent
We benchmark a label-free autonomous agent against uniform random sampling on the task of
discovering distinct topological phases of the long-range Kitaev chain in symmetry class BDI.
Across a $10\times4$ grid of phase counts $d\in[20,200]$ and query budgets
$10^4$--$3\times10^5$, the agent discovers 5--25 more phases than uniform sampling at equal
budget, and the margin is \emph{independent of budget} across a 30-fold range. The agent at
$10^4$ solves finds more phases than uniform sampling at $3\times10^5$ --- a measured
$30\times$ query advantage. Measuring the induced probability of each winding sector with
$4\times10^6$ uniform draws, we find the agent reaches sectors that those draws never hit
($p<2.5\times10^{-7}$): 8 such sectors at $d=140$ from a random start, 21 from the trivial
limit. We separate the rotation-invariant part of the advantage from the part that comes from
knowing where the trivial limit is, and we list the claims from earlier drafts that these
measurements retract.
\end{abstract}

\section{Setup}

After a Jordan--Wigner transformation the generalized cluster chain is the long-range Kitaev
chain in class BDI,
\[
H=i\sum_{j}\sum_{\alpha=0}^{d-1} c_\alpha\, a_j b_{j+\alpha},
\qquad c=(t_0,-t_1,\dots,-t_{d-1}),\qquad \|t\|=1 .
\]
Class BDI is $\mathbb{Z}$-classified in one dimension. The invariant is the winding number of
$f(k)=\sum_\alpha c_\alpha e^{ik\alpha}$, equal to the number of roots of the associated
polynomial inside the unit disc; it takes $d$ values, with $e_0\leftrightarrow\omega=0$
(trivial paramagnet) and $e_{d-1}\leftrightarrow\omega=d-1$. \textbf{All labels in this report
are computed by exact companion-matrix root counting}, not by phase unwrapping on a $k$-grid,
so they carry no lattice-resolution error.

In the Neveu--Schwarz sector $M(t)$ is a twisted circulant, so the ground-state covariance ---
the polar factor of $M$ --- is closed form,
\[
g(r)=\frac{1}{L}\sum_k \frac{f(k)}{|f(k)|}\,e^{-ikr},
\]
verified against explicit polar decomposition to $6\times10^{-15}$ with $\|g\|=1$ exactly.
A ground state at $L=1000$ costs $2.2$\,ms. Throughout: $L=500$, lazy factor $k=30$
(the ground state is re-solved, and one query spent, every 30th parameter step; the gradient is
reused in between, and coverage is counted only at solved points).

\paragraph{Baseline.} Uniform sampling draws $t$ from the rotation-invariant measure
(Gaussian, normalised to the unit sphere) and labels each draw. This choice matters: sampling
along coordinate axes or from heavy-tailed measures presupposes the coupling basis, and the
advantage measured against those baselines does not survive a hidden orthogonal rotation.

\paragraph{Accounting.} One query $=$ one ground-state solve, for both methods. Agent points
are the median over 5 independent seeds (4 in the start-decomposition runs); uniform points are
the median over 5 independent realisations. Bands are min--max, never standard errors.

\section{The grid}

\begin{table}[h]\centering\small
\caption{Median phases discovered, agent ($e_0$ start) / uniform, and the margin. Bold where
the agent's \emph{lowest} seed exceeds uniform's \emph{highest} realisation.}
\begin{tabular}{r rrrr}
\toprule
$d$ & $10^4$ & $3\times10^4$ & $10^5$ & $3\times10^5$\\
\midrule
20 & 19/18 \ (+1) & 20/19 \ (+1) & 20/20 \ (0) & 20/20 \ (0)\\
40 & 32/27 \ \textbf{(+5)} & 34/28 \ \textbf{(+6)} & 35/30 \ \textbf{(+5)} & 37/31 \ \textbf{(+6)}\\
60 & 43/32 \ \textbf{(+11)} & 45/35 \ \textbf{(+10)} & 47/37 \ \textbf{(+10)} & 49/38 \ \textbf{(+11)}\\
80 & 53/37 \ \textbf{(+16)} & 55/40 \ \textbf{(+15)} & 58/43 \ \textbf{(+15)} & 60/45 \ \textbf{(+15)}\\
100 & 59/42 \ \textbf{(+17)} & 62/44 \ \textbf{(+18)} & 67/49 \ \textbf{(+18)} & 68/51 \ \textbf{(+17)}\\
120 & 64/45 \ \textbf{(+19)} & 67/47 \ \textbf{(+20)} & 70/52 \ \textbf{(+18)} & 75/55 \ \textbf{(+20)}\\
140 & 67/47 \ \textbf{(+20)} & 71/52 \ \textbf{(+19)} & 78/56 \ \textbf{(+22)} & 80/59 \ \textbf{(+21)}\\
160 & 74/52 \ \textbf{(+22)} & 77/57 \ \textbf{(+20)} & 80/59 \ \textbf{(+21)} & 88/63 \ \textbf{(+25)}\\
180 & 72/54 \ \textbf{(+18)} & 77/57 \ \textbf{(+20)} & 86/63 \ \textbf{(+23)} & 91/66 \ \textbf{(+25)}\\
200 & 78/59 \ \textbf{(+19)} & 83/62 \ \textbf{(+21)} & 87/67 \ \textbf{(+20)} & 96/71 \ \textbf{(+25)}\\
\bottomrule
\end{tabular}
\end{table}

In 36 of the 40 cells the agent's lowest seed exceeds uniform's highest realisation. The four
exceptions are all $d=20$, where both methods saturate at $20/20$ from $3\times10^4$ upwards;
at $10^4$ the agent loses the worst case there ($18$ vs $20$), because a trajectory still has
to travel while $10^4$ uniform draws already cover a 20-phase space.

\paragraph{The margin does not come from the budget.} Over a 30-fold budget range the margin at
fixed $d$ is flat to within seed noise ($d=200$: $+19,+21,+20,+25$). This closes the objection
that the advantage is bought with compute: it is a difference in sampling efficiency, not in
saturation.

\section{Sampling efficiency}

\begin{table}[h]\centering\small
\caption{The agent at $10^4$ solves against uniform sampling at $3\times10^5$ --- a $30\times$
budget difference in uniform's favour.}
\begin{tabular}{r rrr}
\toprule
$d$ & agent @ $10^4$ & uniform @ $3\times10^5$ & surplus\\
\midrule
60 & 43 & 38 & \ok 5\\
100 & 59 & 51 & \ok 8\\
140 & 67 & 59 & \ok 8\\
200 & 78 & 71 & \ok 7\\
\bottomrule
\end{tabular}
\end{table}

$30\times$ is a \emph{measured lower bound}: $3\times10^5$ is the largest budget we ran.
Uniform coverage grows logarithmically in budget (four-point fit, $R^2=0.95$--$0.99$, slope
$0.6$--$3.7$ phases per $e$-fold); extrapolating that fit, uniform would need $10^6$--$10^7$
draws to match the agent at $10^4$. The extrapolation assumes the logarithmic law survives past
$3\times10^5$ and ignores the hard ceiling at $d$; only the $30\times$ is measured, and only the
$30\times$ is claimed.

\begin{figure}[h]\centering
\includegraphics[width=\textwidth]{fig_main.pdf}
\caption{\textbf{(a)} Coverage across the grid; agent solid, uniform dashed, colour encodes
budget. \textbf{(b)} Coverage against budget at $d=200$: the agent's $10^4$ point lies above
uniform's $3\times10^5$ point. \textbf{(c)} Decomposition of the margin by starting point.}
\end{figure}

\section{Where the margin comes from}

An earlier draft concluded that ``the advantage is the corner start, not the exploration'',
generalising from $d=60$ alone. That generalisation was wrong.

\begin{table}[h]\centering\small
\caption{Start decomposition at $10^5$ solves, median over 4 seeds. The basis-free component
is the margin a \emph{random} start achieves; it grows with $d$ and is not capped.}
\begin{tabular}{r rrr rr}
\toprule
$d$ & agent, $e_0$ & agent, random & uniform & basis-free & extra from $e_0$\\
\midrule
60 & 47.0 & 40.5 & 37 & \textbf{+3.5} & +6.5\\
100 & 65.5 & 52.5 & 49 & \textbf{+3.5} & +13.0\\
140 & 77.5 & 62.5 & 56 & \textbf{+6.5} & +15.0\\
200 & 86.0 & 78.0 & 67 & \textbf{+11.0} & +8.0\\
\bottomrule
\end{tabular}
\end{table}

At $d=200$ the basis-free component exceeds the $e_0$ bonus. Worst-case separation for the
random start holds at $d\geq100$ ($d=200$: worst agent seed 73 $>$ best uniform realisation 71;
$d=140$: $60>58$; $d=100$: $51>49$); $d=60$ is a boundary tie ($38=38$) --- exactly the cell the
earlier draft generalised from.

The margin is therefore a sum of two terms: a rotation-invariant exploration gain that grows
with $d$, and an additional gain from starting at the trivial limit. The second is legitimate
physics --- any experiment starts from the decoupled limit --- but it is basis-dependent
knowledge and is labelled as such throughout.

The two starts are redundant but not identical. At $d=200$ the $e_0$ union over 4 seeds is 103
sectors, the random-start union 93, their intersection 86 and their union 110 (Jaccard $0.78$):
$e_0$ extends into the low tail, random starts into the high tail.

\section{The agent reaches rare sectors}

Uniform sampling induces a sharply peaked measure $p_\omega$ on the $d$ winding sectors. We
measured it with $4\times10^6$ draws per $d$, giving a detection floor of $2.5\times10^{-7}$.

\begin{table}[h]\centering\small
\caption{Sectors reached. ``never seen'' counts sectors the agent reached that $4\times10^6$
uniform draws never hit ($p<2.5\times10^{-7}$); $\omega=0$ is excluded from the $e_0$ column,
since it is the starting point rather than a discovery.}
\begin{tabular}{r rr rr rr}
\toprule
& \multicolumn{2}{c}{uniform} & \multicolumn{2}{c}{agent, $e_0$ @ $10^5$} & \multicolumn{2}{c}{agent, random @ $10^5$}\\
\cmidrule(lr){2-3}\cmidrule(lr){4-5}\cmidrule(lr){6-7}
$d$ & @ $10^5$ & @ $4\times10^6$ & reached & never seen & reached & never seen\\
\midrule
60 & 37 & 43 & 51 & 7 & 45 & 3\\
100 & 49 & 56 & 74 & 18 & 58 & 5\\
140 & 56 & 66 & 88 & 21 & 74 & 8\\
\bottomrule
\end{tabular}
\end{table}

\begin{figure}[h]\centering
\includegraphics[width=\textwidth]{fig_rarity.pdf}
\caption{The measured sector distribution and the reach of each method. Filled dots mark
sectors never seen in $4\times10^6$ uniform draws.}
\end{figure}

\paragraph{Two checks.} \emph{Are the rare labels artefacts?} No. Labels come from exact root
counting, and the structure is coherent: uniform's support is a near-contiguous band and every
extra sector the agent reaches sits on that band's shoulders. Label noise would produce
scattered outliers, not a clean extension. \emph{Is each rare find a single lucky seed?} For
the $e_0$ start, no: 14 of the 21 extra sectors at $d=140$ were found independently by at least
2 of 4 seeds (7 of 7 at $d=60$). For the random start the individual finds are lottery-like
(2 of 8 repeated), though the fact that \emph{some} tail sector is reached is stable.

\paragraph{How rare.} The measured statement is $p<2.5\times10^{-7}$. A quadratic fit to
$\log p_\omega$ (residual rms $0.03$--$0.04$ in log space over the measured range) extrapolates
to far smaller values in the extreme tail, but that is an extrapolation over ten orders of
magnitude with no reason for the large-deviation rate function to remain Gaussian. It is not
claimed here.

\paragraph{Mechanism, stated plainly.} The agent does not snipe isolated exotic sectors. Both
methods cover a contiguous band around the mode; the agent's band is wider, and the sectors at
its edge are quantitatively rare. Both halves of that sentence are needed.

\section{Limitations and retractions}

\begin{itemize}\itemsep2pt
\item \textbf{Retracted:} ``the advantage is the corner start, not the exploration.''
Superseded by the start decomposition above.
\item \textbf{Retracted:} the earlier crossover result. That arm used a memory built on the
accumulated phase $\phi$, which is topologically complete ($\phi_{\text{total}}=2\pi\omega$)
and therefore label-equivalent.
\item \textbf{Not isolated:} ``an underfit autoencoder explores better.'' $L$, bootstrap count,
latent dimension and autoencoder epochs have each been ruled out as the cause of the observed
difference; hidden width $32$ vs $64$ remains untested. Do not cite this claim.
\item \textbf{Scope:} everything here exploits exact solvability of the free-fermion chain.
The interacting ($\kappa$) case needs differentiable DMRG, which is established in the
literature but not done here.
\item \textbf{Scope:} $d=20$ carries no signal --- both methods saturate.
\item \textbf{Open:} $d=200$ sector statistics were still running when this was compiled; the
rarity table covers $d\in\{60,100,140\}$.
\end{itemize}

\section{Reproduction}

\begin{tabular}{@{}ll@{}}
\toprule
\texttt{final\_res/f\_*.json} & agent runs, all $d$ and budgets\\
\texttt{unif\_res/u\_*.json}, \texttt{u1\_*.json} & uniform baselines (consolidated in \texttt{unif\_all.json})\\
\texttt{which\_res/w\_*.json} & start-type runs, with full sector sets\\
\texttt{hist\_res/h\_*.json} & sector histograms, $4\times10^6$ draws per $d$\\
\texttt{plot\_main.py}, \texttt{plot\_rarity.py} & the two figures\\
\texttt{REPORT.md} \S23 & full narrative, including superseded sections\\
\bottomrule
\end{tabular}

\end{document}
